% !TeX root = ../papers/nc-p3.tex
\section{平方座標の三つの境界による負付値縦切片の完全判定}
\label{sec:rank-fifteen-negative-x-fan}
\begin{theorem}\label{thm:rank-fifteen-negative-x-fan}
標数零の通常な階数15窓で、3の上のある付値に対して $\nu(x)<0$ なら、
任意の $y\ne0$ で標準Koszul左右AS正則である。各始点は四中心四次式上で自由階数256である。
\end{theorem}
\begin{proof}
$\nu(y)\le2\nu(x)$ は既知なので、前節の $z=1/B^2$, $u=rB$ がともに正付値の場合を考える。
$(\nu(u),\nu(z))$ 平面で $z=mu^d$、$d=1,2,3$ の三つの境界を用いる。
$d=3$ の表示は $m$ を整とすると、$\nu(z)\ge3\nu(u)$ の全範囲を含む。

\path{computations/verify_rank_fifteen_negative_x_fan.py} は、まず $m$ が一般元のとき
$\mathbb Q(i)(m)[[u]]$ で消去する。係数を実部と虚部に分け、それぞれ
$\mathbb Z[m]$ 上の共通分母で保持する。主元は3と $m$ の奇数次既約剰余式の上で単元に取る。
この剰余体では $-1$ は非平方なので、実部・虚部の少なくとも一方が単元ならよい。
逆元を取るときに使うノルム（実または純虚のときはその非零な座標）の
定数係数の分子・分母を記録し、$u$ 除算では両定数係数の標数零での消滅を検査する。
前節の根の定数項を使うため、初期分母は2の冪しか必要としない。

$d=1$ では十九桁を用い、未切断経路、四次中心、五次の積、六次から十三次の順の損失は
$(0,1,1,2,2)$、$(0,1,2,2)$、二桁、$(3,1,1,1,1,0,1,0)$ である。
$d=2,3$ では共に二十一桁で、これらは
$(0,1,1,3,2)$、$(0,1,2,2)$、二桁、$(4,1,1,1,1,0,1,0)$ となる。
主元多項式の剰余の平方自由な積は順に $m(m^4+1)$、$m(m+1)$、$1$ である。
第一の辺では $m^4+1=(m^2+m+2)(m^2+2m+2)$ の四根を、
第二の辺では $m=-1$ を検査する。第三の辺は $m=0$ を含め例外を持たない。
記録した多項式が消えない剰余では同じ計算が使える。
例外剰余ではその根を含む有限体で計算し、中心の自由列とそれを指定する語基底、
全308行列でのSmith指数列を重複度まで一般の辺と照合する。
補題\ref{lem:rank-fifteen-smith-transfer}を $m$ に関する局所環と一様化変数 $u$ に適用すれば、
この比較で一致した剰余の上の全ての持上げでも消去できる。
$d=1,2$ では $m=0$ は開いた隣接領域に属し、境界で検査するのは単元だけである。

残る三錐の辺の原始ベクトルは $(1,0),(1,1),(1,2),(1,3)$ で、
隣り合う二辺の行列式は1である。$(1,0)$ の辺では $z$ は一般元であり、
$h^2=z$ を添加すれば $r=hu$ が $u$ の単元倍である。
第\ref{sec:rank-fifteen-unit-pole-depth}節の一般円錐曲線で $v=A/B$,
$(Q_+/u,Q_-/u)|_{u=0}=(1+v,1-v)$ となり、同じ符号の規格化を与える。
この既知の辺と三つの一般境界を順に使う。

各錐の対角線 $(u,z)=(t^2,t),(t^2,t^3),(t^2,t^5)$ で、剰余体 $\mathbb F_9$ 上それぞれ
35、39、41桁を用いる。いずれも必要な中心の自由列とそれを指定する語基底が両辺で一致し、
全308行列について、昇順の指数の項別和が対角線の全指数に一致する。
補題\ref{lem:rank-fifteen-two-parameter-smith}は $i$ を添加した3上の非分岐局所環でも同じ証明で成り立つ。
各錐の単項式座標で未切断経路、四次中心、左右の積と中心商へ順次適用する。
十三次での消滅と全始点の上界から、従来の全次数判定で結論が従う。
\end{proof}

\begin{corollary}
$\nu(a)<0$ となる3の上の付値がある縦切片 $x=a$ では、全非零水準の添字保持同型類が
系\ref{cor:rank-fifteen-unit-vertical-conics}の $T$ の式で分類される。
分岐水準 $a^3/27,(1-a)^2/4$ で各二類、その他で各四類となり、最小周期は全て2である。
さらに任意の有理数 $x$ と非零の代数的数 $y$ で、$\mathbb Q(y)$ が
3上に剰余次数4以下の素点を持てば、全ての通常窓が上の正則性と自由性を持つ。
特に全ての有理座標対 $(x,y)$ が認証される。
\end{corollary}
\begin{proof}
通常表示外の四点の横座標 $3,3/4$ および $4a^2-12a+7=0$ の二根は全て3整である。
$(1-a)^2-4a^3/27$ の零点も $3,3/4$ なので、定理と既知の円錐曲線表示が縦切片の全点で使える。
数体については、$x$ が3整なら第\ref{sec:rank-fifteen-negative-unit-levels}節の系、
負付値なら上の定理による。
\end{proof}
